\itemhead{1.1 例题 · 例题1（教材p.3）}{XJ-B1-C01-EX-P003-Y0620}
\label{XJ-B1-C01-EX-P003-Y0620}
{\footnotesize 湘教版教材p.3 / PDF第10页 / 1.1 例题 / 例题1}
\field{原题}{设 $L(A,B)$ 表示直线 $AB$ 上全体点组成的集合，则 $P\in L(A,B)$ 的含义是什么？}
\field{参考答案}{$P$ 是直线 $AB$ 上的一个点。}
{\small \field{关联知识}{S01 集合、元素与常用数集（教材帮p.1、2 / PDF10、11）}
\field{方法与易错点}{先辨认研究对象和数集范围，再代入定义判断；含参集合逐个来源分类。 易错：把自然数集误作正整数集；混淆元素和单元素集合；候选参数造成元素重复。}
\field{答案依据}{例题原解位置：湘教版教材p.3及题库原图；独立编写可编辑参考解答。答案册此部分主要收练习、习题，不伪造例题答案册位置。}}
